\documentclass[a4paper,10pt]{paper}
%DIF LATEXDIFF DIFFERENCE FILE
%DIF DEL problem2.tex           Fri Oct 24 03:07:42 2025
%DIF ADD revised/problem2.tex   Thu Sep  3 07:05:52 2026

\usepackage{mycommands}
%DIF 4d4
%DIF < \usepackage{amsmath}
%DIF -------

\title{\DIFaddbegin \DIFadd{Seismology }\DIFaddend Problem Set 2}
\author{David Al-Attar \\ Michaelmas Term, \the\year}
%DIF PREAMBLE EXTENSION ADDED BY LATEXDIFF
%DIF UNDERLINE PREAMBLE %DIF PREAMBLE
\RequirePackage[normalem]{ulem} %DIF PREAMBLE
\RequirePackage{color}\definecolor{RED}{rgb}{1,0,0}\definecolor{BLUE}{rgb}{0,0,1} %DIF PREAMBLE
\providecommand{\DIFadd}[1]{{\protect\color{blue}\uwave{#1}}} %DIF PREAMBLE
\providecommand{\DIFdel}[1]{{\protect\color{red}\sout{#1}}}                      %DIF PREAMBLE
%DIF SAFE PREAMBLE %DIF PREAMBLE
\providecommand{\DIFaddbegin}{} %DIF PREAMBLE
\providecommand{\DIFaddend}{} %DIF PREAMBLE
\providecommand{\DIFdelbegin}{} %DIF PREAMBLE
\providecommand{\DIFdelend}{} %DIF PREAMBLE
\providecommand{\DIFmodbegin}{} %DIF PREAMBLE
\providecommand{\DIFmodend}{} %DIF PREAMBLE
%DIF FLOATSAFE PREAMBLE %DIF PREAMBLE
\providecommand{\DIFaddFL}[1]{\DIFadd{#1}} %DIF PREAMBLE
\providecommand{\DIFdelFL}[1]{\DIFdel{#1}} %DIF PREAMBLE
\providecommand{\DIFaddbeginFL}{} %DIF PREAMBLE
\providecommand{\DIFaddendFL}{} %DIF PREAMBLE
\providecommand{\DIFdelbeginFL}{} %DIF PREAMBLE
\providecommand{\DIFdelendFL}{} %DIF PREAMBLE
%DIF COLORLISTINGS PREAMBLE %DIF PREAMBLE
\RequirePackage{listings} %DIF PREAMBLE
\RequirePackage{color} %DIF PREAMBLE
\lstdefinelanguage{DIFcode}{ %DIF PREAMBLE
%DIF DIFCODE_UNDERLINE %DIF PREAMBLE
  moredelim=[il][\color{red}\sout]{\%DIF\ <\ }, %DIF PREAMBLE
  moredelim=[il][\color{blue}\uwave]{\%DIF\ >\ } %DIF PREAMBLE
} %DIF PREAMBLE
\lstdefinestyle{DIFverbatimstyle}{ %DIF PREAMBLE
	language=DIFcode, %DIF PREAMBLE
	basicstyle=\ttfamily, %DIF PREAMBLE
	columns=fullflexible, %DIF PREAMBLE
	keepspaces=true %DIF PREAMBLE
} %DIF PREAMBLE
\lstnewenvironment{DIFverbatim}{\lstset{style=DIFverbatimstyle}}{} %DIF PREAMBLE
\lstnewenvironment{DIFverbatim*}{\lstset{style=DIFverbatimstyle,showspaces=true}}{} %DIF PREAMBLE
%DIF END PREAMBLE EXTENSION ADDED BY LATEXDIFF

\begin{document}

\maketitle

\begin{enumerate}

\item Consider the regularised least squares solution to a linear(ised) inverse problem obtained by minimising the function
\begin{equation*}
    J(\mathbf{m})=\frac{1}{2}(\mathbf{A}\mathbf{m}-\mathbf{d})^{T}\mathbf{C}^{-1}(\mathbf{A}\mathbf{m}-\mathbf{d})+\frac{\lambda}{2}\mathbf{m}^{T}\mathbf{B}\mathbf{m}.
\end{equation*}
Suppose that the data \DIFdelbegin \DIFdel{takes the form $\mathbf{d}=\mathbf{A}\mathbf{m}_{in}$ with $\mathbf{m}_{in}$ }\DIFdelend \DIFaddbegin \DIFadd{take the form $\mathbf{d}=\mathbf{A}\mathbf{m}_{\mathrm{in}}$ with $\mathbf{m}_{\mathrm{in}}$ }\DIFaddend a specified input model. Show that the resulting least squares solution, \DIFdelbegin \DIFdel{$\mathbf{m}_{out}$}\DIFdelend \DIFaddbegin \DIFadd{$\mathbf{m}_{\mathrm{out}}$}\DIFaddend , takes the form
\DIFdelbegin \begin{displaymath}
    \DIFdel{\mathbf{m}_{out}=\mathbf{R}\mathbf{m}_{in},
}\end{displaymath}%DIFAUXCMD
\DIFdelend \DIFaddbegin \begin{equation*}
    \DIFadd{\mathbf{m}_{\mathrm{out}}=\mathbf{R}\mathbf{m}_{\mathrm{in}},
}\end{equation*}\DIFaddend 
where $\mathbf{R}$ is a matrix to be determined. Discuss the significance of this so-called resolution matrix, explaining why, in particular, it would be desirable for its value to be close to the identity matrix.

\item

Consider a linear inverse problem of the standard form
\begin{equation*}
    \mathbf{d} = \mathbf{A}\mathbf{m}+\mathbf{e},
\end{equation*}
with $\mathbf{d}$ an $n$-dimensional data vector, $\mathbf{m}$ an $m$-dimensional model vector, $\mathbf{A}$ an $n\times m$ matrix, and $\mathbf{e}$ \DIFdelbegin \DIFdel{and }\DIFdelend \DIFaddbegin \DIFadd{an }\DIFaddend $n$-dimensional vector of random errors. We assume $m > n$, which
is to say that there are more unknown model parameters than available data.
We further suppose that the data \DIFdelbegin \DIFdel{has }\DIFdelend \DIFaddbegin \DIFadd{have }\DIFaddend been processed such that the components of the random errors are uncorrelated and have unit standard deviations (this corresponds to the data covariance, $\mathbf{C}$, being the $n\times n$ identity matrix).



\begin{enumerate}

    \item A model $\mathbf{m}$ is said to be compatible with the data if
\begin{equation*}
    \|\mathbf{A}\mathbf{m}-\mathbf{d}\|^{2} \le \chi^{2}_{c},
\end{equation*}
where $\chi^{2}_{c}$ is a critical value for the so-called chi-squared statistic. What type of geometric region of the \DIFdelbegin \DIFdel{model-space }\DIFdelend \DIFaddbegin \DIFadd{model space }\DIFaddend is defined by this inequality?

\item Supposing that \DIFdelbegin \DIFdel{$ \|\mathbf{d}\|^{2} > \chi^{2}_{c}$}\DIFdelend \DIFaddbegin \DIFadd{$\|\mathbf{d}\|^{2} > \chi^{2}_{c}$}\DIFaddend , the region of the model space defined by the inequality in (a) does not contain the origin. In such
cases, the closest point within this region to the origin must lie on the \DIFdelbegin \DIFdel{boundary-set }\DIFdelend \DIFaddbegin \DIFadd{boundary set }\DIFaddend defined by
\begin{equation*}
    \|\mathbf{A}\mathbf{m}-\mathbf{d}\|^{2} = \chi^{2}_{c}.
\end{equation*}
Using the method of Lagrange multipliers, we can determine this \emph{minimum norm solution} by considering the Lagrangian
\begin{equation*}
    \mathcal{L}(\mathbf{m},\mu)=\|\mathbf{m}\|^{2}+\mu(\|\mathbf{A}\mathbf{m}-\mathbf{d}\|^{2}-\chi^{2}_{c}),
\end{equation*}
where $\mu$ is a Lagrange multiplier. Derive the equations for $\mathbf{m}$ and $\mu$ that must hold at the minimum norm solution.
Discuss the relationship between this approach and regularised least squares.
\end{enumerate}


\item Consider the deformation of an isotropic elastic body due to an applied surface load, $\sigma$. The relevant linearised equations of motion are
\DIFdelbegin \begin{eqnarray*}
    \DIFdel{-\frac{\partial T_{ij}}{\partial x_{j}}=0; }\\
    \DIFdel{T_{ij}=\lambda\frac{\partial u_{k}}{\partial x_{k}}\delta_{ij}+\mu\left(\frac{\partial u_{i}}{\partial x_{j}}+\frac{\partial u_{j}}{\partial x_{i}}\right),
}\end{eqnarray*}%DIFAUXCMD
\DIFdelend \DIFaddbegin \begin{gather*}
    \DIFadd{-\frac{\partial T_{ij}}{\partial x_{j}}=0, }\\
    \DIFadd{T_{ij}=\lambda\frac{\partial u_{k}}{\partial x_{k}}\delta_{ij}+\mu\left(\frac{\partial u_{i}}{\partial x_{j}}+\frac{\partial u_{j}}{\partial x_{i}}\right),
}\end{gather*}\DIFaddend 
within the reference body $M$, and subject to boundary conditions
\begin{equation*}
    T_{ij}\hat{n}_{j}=\sigma g_{i}
\end{equation*}
on its boundary, $\partial M$. Here $g_{i}$ is the acceleration due to gravity prior to the deformation. Note that self-gravitation is being ignored within these equations.

Suppose that observations, \DIFdelbegin \DIFdel{$\mathbf{u}_{i}^{obs}$}\DIFdelend \DIFaddbegin \DIFadd{$\mathbf{u}_{i}^{\mathrm{obs}}$}\DIFaddend , of the displacement vector have been made at surface locations $\mathbf{x}_i$ for \DIFdelbegin \DIFdel{$i=1,...,m$}\DIFdelend \DIFaddbegin \DIFadd{$i=1,\dots,m$}\DIFaddend . An inverse problem could then be formulated to estimate the Lamé parameters, $\lambda$ and $\mu$, as well as the load, $\sigma$. To this end, we can seek to minimise the misfit
\DIFdelbegin \begin{displaymath}
    \DIFdel{J=\frac{1}{2}\sum_{i=1}^{m}\frac{1}{\sigma_{i}^2}||\mathbf{u}(\mathbf{x}_{i})-\mathbf{u}_{i}^{obs}||^{2}
}\end{displaymath}%DIFAUXCMD
\DIFdelend \DIFaddbegin \begin{equation*}
    \DIFadd{J=\frac{1}{2}\sum_{i=1}^{m}\frac{1}{\varepsilon_{i}^{2}}\|\mathbf{u}(\mathbf{x}_{i})-\mathbf{u}_{i}^{\mathrm{obs}}\|^{2},
}\end{equation*}\DIFaddend 
where \DIFdelbegin \DIFdel{$\sigma_{i}$ }\DIFdelend \DIFaddbegin \DIFadd{$\varepsilon_{i}$ }\DIFaddend is the standard deviation for the $i$\DIFdelbegin \DIFdel{-th }\DIFdelend \DIFaddbegin \DIFadd{th }\DIFaddend observation. Adapt the approach within Lecture\DIFaddbegin \DIFadd{~}\DIFaddend 19 to derive sensitivity kernels for $J$ with respect to $\lambda$, $\mu$, and $\sigma$, showing how they can be calculated through one solution of the forward problem and one of an appropriate adjoint problem.





\item Recall that the Lagrangian for a self-gravitating elastic body is given by
\DIFdelbegin \begin{displaymath}
    \DIFdel{\mathcal{L}(\mathbf{x},t,\boldsymbol{\varphi},\mathbf{v},\mathbf{F})=\frac{1}{2}\rho(\mathbf{x})||\mathbf{v}||^{2}-W(\mathbf{x},t,\mathbf{F})-\frac{1}{2}\rho(\mathbf{x})\zeta(\mathbf{x},t),
}\end{displaymath}%DIFAUXCMD
\DIFdelend \DIFaddbegin \begin{equation*}
    \DIFadd{\mathcal{L}(\mathbf{x},t,\bphi,\mathbf{v},\mathbf{F})=\frac{1}{2}\rho(\mathbf{x})\|\mathbf{v}\|^{2}-W(\mathbf{x},t,\mathbf{F})-\frac{1}{2}\rho(\mathbf{x})\zeta(\mathbf{x},t),
}\end{equation*}\DIFaddend 
where $\zeta$ is the referential gravitational potential
\begin{equation*}
    \zeta(\mathbf{x},t)=-G\int_{M}\frac{\rho(\mathbf{x}')}{\{[\varphi_{i}(\mathbf{x},t)-\varphi_{i}(\mathbf{x}',t)][\varphi_{i}(\mathbf{x},t)-\varphi_{i}(\mathbf{x}',t)]\}^{\frac{1}{2}}}\dd^{3}\mathbf{x}'.
\end{equation*}
By calculating the appropriate functional derivative, show that
\DIFdelbegin \begin{displaymath}
    \DIFdel{\frac{\delta\mathcal{L}}{\delta\varphi_{i}}=\rho\gamma_{i}
}\end{displaymath}%DIFAUXCMD
\DIFdelend \DIFaddbegin \begin{equation*}
    \DIFadd{\frac{\delta\mathcal{L}}{\delta\varphi_{i}}=\rho\gamma_{i},
}\end{equation*}\DIFaddend 
where $\gamma_{i}$ is the referential gravitational acceleration.

\item Consider a planet at equilibrium occupying the volume $M$ and with density $\rho$. Its gravitational potential, $\phi$, is a solution of Poisson's equation
\begin{equation*}
    \nabla^{2}\phi=4\pi G\rho,
\end{equation*}
where it is understood that $\rho$ vanishes outside of $M$. Across the surface $\partial M$ we then have continuity conditions on $\phi$ and its normal derivative, and finally require that $\phi$ tend to zero at infinity. Suppose that the planet is only slightly aspherical, with boundary
\DIFdelbegin \begin{displaymath}
    \DIFdel{r=b+sh_{1}(\theta,\varphi)+\cdot\cdot\cdot,
}\end{displaymath}%DIFAUXCMD
\DIFdelend \DIFaddbegin \begin{equation*}
    \DIFadd{r=b+sh_{1}(\theta,\varphi)+\cdots,
}\end{equation*}\DIFaddend 
in spherical polar co-ordinates, and its density decomposed as
\DIFdelbegin \begin{displaymath}
    \DIFdel{\rho=\rho_{0}+s \rho_{1}+\cdot\cdot\cdot,
}\end{displaymath}%DIFAUXCMD
\DIFdelend \DIFaddbegin \begin{equation*}
    \DIFadd{\rho=\rho_{0}+s\rho_{1}+\cdots,
}\end{equation*}\DIFaddend 
with $\rho_{0}$ spherically symmetric. Here $s$ is a perturbation parameter, and the dots denote higher-order terms that are to be neglected. Similarly writing the potential in the form \DIFdelbegin \DIFdel{$\phi=\phi_{0}+s \phi_{1}+\cdot\cdot\cdot$ }\DIFdelend \DIFaddbegin \DIFadd{$\phi=\phi_{0}+s\phi_{1}+\cdots$, }\DIFaddend show that across the reference surface $r=b$ we require continuity of both $\phi_{1}$ and
\begin{equation*}
    \frac{\partial\phi_{1}}{\partial r}+4\pi G\rho_{0}h_{1}.
\end{equation*}

\item Using results from Lecture\DIFaddbegin \DIFadd{~}\DIFaddend 22 along with residue calculus (which you don't need to know for this course), it can be shown that the time-domain solution of the elastodynamic equations in a non-rotating earth model can be expressed as the following eigenfunction expansion
\begin{equation*}
    \mathbf{u}(\mathbf{x},t)=\sum_{km}\left\{\int_{0}^{t}\frac{\sin[\omega_{k}(t-t')]}{\omega_{k}}\int_{M}\overline{S}_{ij}(\mathbf{x}',t')\frac{\partial|km\rangle^{*}}{\partial x'_{j}}(\mathbf{x}')\dd^{3}\mathbf{x}'\dd t'\right\}|km\rangle(\mathbf{x}).
\end{equation*}
Simplify this expression in the case of a point source for which
\begin{equation*}
    \overline{S}_{ij}(\mathbf{x},t)=M_{ij}\delta(\mathbf{x}-\mathbf{x}_{s})H(t-t_{s}),
\end{equation*}
with $M_{ij}$ known as the moment tensor, $\mathbf{x}_{s}$ the source location, and $t_{s}$ the source time. Describe briefly and qualitatively how an inverse problem might be formulated to estimate the ten source parameters $(M_{ij},\mathbf{x}_{s},t_{s})$ from recorded seismograms.

\item Consider the eigenvalue problem for a rotating earth model
\DIFdelbegin \begin{displaymath}
    \DIFdel{-\omega^{2}\langle \mathbf{w}|\mathbf{P}|\mathbf{s}\rangle+\ii\omega\langle \mathbf{w}|\mathbf{W}|\mathbf{s}\rangle+\langle \mathbf{w}|\mathbf{H}|\mathbf{s}\rangle=0
}\end{displaymath}%DIFAUXCMD
\DIFdelend \DIFaddbegin \begin{equation*}
    \DIFadd{-\omega^{2}\langle \mathbf{w}|P|\mathbf{s}\rangle+\ii\omega\langle \mathbf{w}|W|\mathbf{s}\rangle+\langle \mathbf{w}|H|\mathbf{s}\rangle=0,
}\end{equation*}\DIFaddend 
where $\mathbf{s}$ is the eigenfunction, $\omega$ the \DIFdelbegin \DIFdel{eigenvalue}\DIFdelend \DIFaddbegin \DIFadd{eigenfrequency}\DIFaddend , and $\mathbf{w}$ an arbitrary test function. Assuming the eigenfunction normalisation condition
\DIFdelbegin \begin{displaymath}
    \DIFdel{\langle \mathbf{s}|\mathbf{P}|\mathbf{s}\rangle=1,
}\end{displaymath}%DIFAUXCMD
\DIFdelend \DIFaddbegin \begin{equation*}
    \DIFadd{\langle \mathbf{s}|P|\mathbf{s}\rangle=1,
}\end{equation*}\DIFaddend 
show that the eigenfrequency can be written
\DIFdelbegin \begin{displaymath}
    \DIFdel{\omega=\frac{1}{2}\ii\langle \mathbf{s}|\mathbf{W}|\mathbf{s}\rangle\pm\sqrt{\langle \mathbf{s}|\mathbf{H}|\mathbf{s}\rangle-\frac{1}{4}\langle \mathbf{s}|\mathbf{W}|\mathbf{s}\rangle^{2}}.
}\end{displaymath}%DIFAUXCMD
\DIFdelend \DIFaddbegin \begin{equation*}
    \DIFadd{\omega=\frac{1}{2}\ii\langle \mathbf{s}|W|\mathbf{s}\rangle\pm\sqrt{\langle \mathbf{s}|H|\mathbf{s}\rangle-\frac{1}{4}\langle \mathbf{s}|W|\mathbf{s}\rangle^{2}}.
}\end{equation*}\DIFaddend 
Use this result to discuss qualitatively the stability of a rotating earth model, including, in particular, the role of the Coriolis and centrifugal forces.

\end{enumerate}

\end{document}
